\section{Introduction}
\label{sec:introduction}

A spanning tree retains connectivity with only $n-1$ edges. A thin
spanning tree must also use a small fraction of every cut, with the
same tree serving all cuts simultaneously. We give a deterministic polynomial-bit
construction at the inverse-connectivity scale.

Throughout, graphs are finite undirected loopless multigraphs, and
every parallel copy counts separately. For $S\subseteq V(G)$, let
$\delta_G(S)$ be the multiset of edges with one endpoint in $S$ and
the other in its complement. The graph is $k$-edge-connected if
$|\delta_G(S)|\ge k$ for every nonempty proper $S$. A spanning tree
$T$ is \emph{$\alpha$-thin} in $G$ if
\[
 |\delta_T(S)|\le\alpha|\delta_G(S)|
 \qquad(\varnothing\ne S\subsetneq V(G)).
\]
We consider two input representations: an explicit list of labelled
edge copies, or a list of endpoint pairs with nonnegative integer
multiplicities written in binary. In the second representation, the
number of edge copies may be exponential in the input length.

\begin{theorem}[Algorithmic strong thin trees]
\label{thm:tt-algorithm}
There is an absolute constant $C$ with the following property. Given a
finite $k$-edge-connected multigraph $G$ on at least one vertex and
an integer $k\ge1$,
a deterministic algorithm returns a spanning tree $T$ satisfying
\begin{equation}
 |\delta_T(S)|\le \frac Ck\,|\delta_G(S)|
       \qquad(\varnothing\ne S\subsetneq V(G))
 \label{eq:ta-thin}
\end{equation}
in time polynomial in the total binary input length, whether edges
are listed explicitly or multiplicities are encoded in binary.
For a one-vertex graph the answer is the empty tree.
\end{theorem}


The universal $C/k$ bound is the strong thin-tree conjecture. The
companion article \emph{The strong thin tree conjecture}
\cite{ThinTreesCompanion} proves the existence statement and supplies
the complete hierarchy geometry used here. Theorem~\ref{thm:tt-algorithm}
gives the positive algorithmic resolution in both finite input
models. Its order in $k$ is necessary: for two vertices joined by
$k$ parallel edges, the unique cut contains exactly one tree edge.

The construction also gives simultaneous cost-and-cut rounding. Given
nonnegative rational edge weights whose mass on every cut is at least one, and any
nonnegative rational edge costs, we find a tree in the positive support
with every cut load and its total cost bounded by constant multiples of
the corresponding fractional quantities
(Corollary~\ref{cor:ta-fractional-rounding}). No metric assumption on
the costs is needed.

\subsection{Context and prior work}

Goddyn's qualitative conjecture asks whether every positive thinness
bound follows from sufficiently large edge connectivity, independently
of graph order \cite[Problem~4]{Goddyn2004}. The strong form specifies
the rate $C/k$; the distinction is made explicitly in
\cite[Conjecture~1.3 and the following discussion]{AnariOveisGharan2015}.
One source of algorithmic interest is the asymmetric traveling salesman
problem, which asks for a minimum-cost tour through the vertices of a
directed metric. Asadpour, Goemans, M\k{a}dry, Oveis Gharan, and Saberi
used maximum-entropy spanning-tree distributions to obtain an
$O(\log n/\log\log n)$ approximation
\cite{AsadpourEtAl2017}. Their rounding method also gives the
constructive $O(\log n/(k\log\log n))$ thin-tree scale; see
\cite[Section~1]{AnariOveisGharan2015}.

For planar graphs, and more generally on each fixed orientable
surface, Oveis Gharan and Saberi obtained constructive inverse-connectivity
bounds \cite[Theorem~5.1]{OveisGharanSaberi2011}.
For general multigraphs, Anari and Oveis Gharan introduced
effective-resistance-reducing flows and locally connected hierarchies,
and proved the existence of
$\operatorname{poly}(\log\log n)/k$-thin trees
\cite[Corollary~1.8]{AnariOveisGharan2015}. Their geometric method is
the principal foundation of the hierarchy argument in the companion
and its extension below. The existence theorem does not itself give
a polynomial-time construction of those trees.

Recent work has obtained inverse-connectivity bounds for structured
families of cuts. Klein and Olver treat every prescribed laminar
family \cite{KleinOlver2023}. Klein, Olver, and Yeoh construct a tree
crossing each cut of size less than $(1+1/40)k$ at most $88$ times
\cite[Theorem~1]{KleinOlverYeoh2026}. Our conclusion controls all cuts
of an arbitrary multigraph. Thin trees remain an independent
structural problem: constant-factor approximation for asymmetric
traveling salesman was already established by Svensson, Tarnawski,
and V\'egh and subsequently improved by Traub and Vygen
\cite{SvenssonTarnawskiVegh2020,TraubVygen2022}.

\subsection{From the geometric existence proof to an algorithm}

High edge connectivity supplies many edge-disjoint spanning trees
by the Nash--Williams--Tutte theorem
\cite{NashWilliams1961,Tutte1961}. A packing bounds the average
tree load on each fixed cut, but does not identify one tree that is
small on every cut. The companion retains a large packing through
successive sparsifications, reducing cut sizes and packing size by
almost the same factor. Two vector representations allow the same
edge subset to preserve cut control and enough rank information to
certify all partition inequalities for tree packing. Iterated matrix
halving with a bounded product of losses was used by Harvey and Olver
to obtain spectrally thin trees under an effective-resistance hypothesis
\cite[full version, Appendix~F]{HarveyOlver2014}. Here the additional
packing representation permits iteration from edge connectivity.

A shortcut matrix is a positive quadratic form that bounds cut sizes;
its inverse assigns each edge a resistance, and edges of low resistance
are the candidates retained for sparsification. Three steps of the
existence argument require new computational work. The shortcut matrix
must be consistent with partitions determined by that same matrix; its
construction by compactness and a fixed point does not provide a uniform
finite search bound. Real whitening may introduce irrational entries,
so a polynomial-bit algorithm must specify its numerical data and control
every precision loss. Finally, the existential matrix partition theorem
of Marcus, Spielman, and Srivastava must be replaced by an algorithm
that finds the signs. We use the public constructive method of Ezeunala
and Jiang, with a quantitative implementation proved below, while
retaining the packing argument.

\paragraph{A larger cone of geometric tests.}
The companion bounds a shortcut matrix on cut vectors. For optimization
we use centered positive semidefinite matrices whose squared distances
satisfy triangle inequalities. This cone has polynomially many displayed
constraints. Enlarging the cone strengthens the required shortcut
inequality, so the companion's cut-only statement cannot simply be
applied. We extend its proof by constructing a doubled distance-profile
polytope in which radii about vertex images are affine and radial motion
has exactly its metric length. The underlying distance-profile isometry
is classical \cite{Frechet1910,Kuratowski1935}; the additional radial
identities make the bag and charging arguments work in the new geometry.

\paragraph{Continuous capacities and finite optimization.}
Instead of selecting edges abruptly at a resistance threshold, we assign
capacities through a continuous ramp. A convex response lowers the
resistance cutoff while maintaining the partition inequalities with a
controlled loss in packing size.
Uncrossing reduces a putative obstruction to a chain of partitions,
where the geometric response applies simultaneously. Exact rational
separation and a quantitative interior ball turn each response into a
finite computation using the ellipsoid method \cite{GLS1988}.
Constructive graphic matroid partition then recovers an integer packing
\cite{Edmonds1968,Terao2023}.

\paragraph{Exact data and certified sparsification.}
A log-determinant normalization, following the variational method behind
radial isotropic position \cite[Section~2.2]{Barthe1998}, gives the second
vector representation. We approximate the stacked vectors rationally,
complete their covariance exactly to the identity, and apply constructive
matrix signing. This step builds on the Kadison--Singer partition theorem
of Marcus, Spielman, and Srivastava
\cite{MarcusSpielmanSrivastava2015} and the public rank-one discrepancy
method of Ezeunala and Jiang \cite{EzeunalaJiang2026}; the exact interface
and the constant used are proved in Appendix~\ref{app:signing}. Rational matrix
certificates imply every cut and packing estimate needed for the next
step. The product of the multiplicative losses remains bounded through
the iteration. A separate truncation and rounding reduction handles
binary multiplicities before any edge expansion: it produces a graph
with polynomially many explicit copies, independently of the numerical
values of the multiplicities and $k$. The final tree lifts to original
edge copies.

\paragraph{Organization.}
Section~\ref{tg:section} proves the metric-cone hierarchy estimate and
its simultaneous response to a partition chain. Section~\ref{sec:ta-algorithm}
constructs a rational shortcut and a large packing of low-resistance edges.
Section~\ref{sec:normalization} builds the rank-preserving frame.
Section~\ref{sec:signing} provides exact rational padding and signing
certificates. Section~\ref{sec:iteration} combines the two representations,
iterates the sparsification, and proves the binary-multiplicity reduction.
Section~\ref{sec:fractional-rounding} records a simultaneous cost-and-cut
rounding consequence for rational fractional inputs.
Appendix~\ref{app:signing} gives the quantitative rational implementation
of the public rank-one signing method.
